\documentclass[a4paper, 14pt]{extarticle}

% Поля
%--------------------------------------
\usepackage{geometry}
\geometry{a4paper,tmargin=2cm,bmargin=2cm,lmargin=3cm,rmargin=1cm}
%--------------------------------------

\usepackage{float}
%Russian-specific packages
%--------------------------------------
\usepackage[T2A]{fontenc}
\usepackage[utf8]{inputenc} 
\usepackage[english, main=russian]{babel}
%--------------------------------------

\usepackage{textcomp}

% Красная строка
%--------------------------------------
\usepackage{indentfirst}               
%--------------------------------------             
\usepackage[most]{tcolorbox}
\tcbuselibrary{listings,breakable}
\newtcblisting{breakablelisting}{
  listing only,
  breakable,
  enhanced,
  frame hidden,
  colback=white,
  listing options={
    mathescape=false,
    texcl=false,
    inputencoding=utf8,
    literate={\$}{{\$}}1
  }
}


%Graphics
%--------------------------------------
\usepackage{graphicx}
\usepackage{wrapfig}
%--------------------------------------
\graphicspath{{./lab5/}}
% Полуторный интервал
%--------------------------------------
\linespread{1.3}                    
%--------------------------------------

%Выравнивание и переносы
%--------------------------------------
% Избавляемся от переполнений
\sloppy
% Запрещаем разрыв страницы после первой строки абзаца
\clubpenalty=10000
% Запрещаем разрыв страницы после последней строки абзаца
\widowpenalty=10000
%--------------------------------------

%Списки
\usepackage{enumitem}

%Подписи
\usepackage{caption} 

%Гиперссылки
\usepackage{hyperref}

\hypersetup {
	unicode=true
}

%Рисунки
%--------------------------------------
\DeclareCaptionLabelSeparator*{emdash}{~--- }
\captionsetup[figure]{labelsep=emdash,font=onehalfspacing,position=bottom}
%--------------------------------------

\usepackage{tempora}

%Листинги
%--------------------------------------
\usepackage{listings}
\lstset{
  basicstyle=\ttfamily\footnotesize, 
  %basicstyle=\footnotesize\AnkaCoder,        % the size of the fonts that are used for the code
  breakatwhitespace=false,         % sets if automatic breaks shoulbd only happen at whitespace
  breaklines=true,                 % sets automatic line breaking
  captionpos=t,                    % sets the caption-position to bottom
  inputencoding=utf8,
  frame=single,                    % adds a frame around the code
  keepspaces=true,                 % keeps spaces in text, useful for keeping indentation of code (possibly needs columns=flexible)
  keywordstyle=\bf,       % keyword style
  numbers=left,                    % where to put the line-numbers; possible values are (none, left, right)
  numbersep=5pt,                   % how far the line-numbers are from the code
  xleftmargin=25pt,
  xrightmargin=25pt,
  showspaces=false,                % show spaces everywhere adding particular underscores; it overrides 'showstringspaces'
  showstringspaces=false,          % underline spaces within strings only
  showtabs=false,                  % show tabs within strings adding particular underscores
  stepnumber=1,                    % the step between two line-numbers. If it's 1, each line will be numbered
  tabsize=2,                       % sets default tabsize to 8 spaces
  title=\lstname                   % show the filename of files included with \lstinputlisting; also try caption instead of title
}
%--------------------------------------

%%% Математические пакеты %%%
%--------------------------------------
\usepackage{amsthm,amsfonts,amsmath,amssymb,amscd}  % Математические дополнения от AMS
\usepackage{mathtools}                              % Добавляет окружение multlined
\usepackage[perpage]{footmisc}
%--------------------------------------

%--------------------------------------
%			НАЧАЛО ДОКУМЕНТА
%--------------------------------------

\begin{document}

%--------------------------------------
%			ТИТУЛЬНЫЙ ЛИСТ
%--------------------------------------
\begin{titlepage}
\thispagestyle{empty}
\newpage


%Шапка титульного листа
%--------------------------------------
\vspace*{-60pt}
\hspace{-65pt}
\begin{minipage}{0.3\textwidth}
\hspace*{-20pt}\centering
\includegraphics[width=\textwidth]{emblem}
\end{minipage}
\begin{minipage}{0.67\textwidth}\small \textbf{
\vspace*{-0.7ex}
\hspace*{-6pt}\centerline{Министерство науки и высшего образования Российской Федерации}
\vspace*{-0.7ex}
\centerline{Федеральное государственное автономное образовательное учреждение }
\vspace*{-0.7ex}
\centerline{высшего образования}
\vspace*{-0.7ex}
\centerline{<<Московский государственный технический университет}
\vspace*{-0.7ex}
\centerline{имени Н.Э. Баумана}
\vspace*{-0.7ex}
\centerline{(национальный исследовательский университет)>>}
\vspace*{-0.7ex}
\centerline{(МГТУ им. Н.Э. Баумана)}}
\end{minipage}
%--------------------------------------

%Полосы
%--------------------------------------
\vspace{-25pt}
\hspace{-35pt}\rule{\textwidth}{2.3pt}

\vspace*{-20.3pt}
\hspace{-35pt}\rule{\textwidth}{0.4pt}
%--------------------------------------

\vspace{1.5ex}
\hspace{-35pt} \noindent \small ФАКУЛЬТЕТ\hspace{80pt} <<Информатика и системы управления>>

\vspace*{-16pt}
\hspace{47pt}\rule{0.83\textwidth}{0.4pt}

\vspace{0.5ex}
\hspace{-35pt} \noindent \small КАФЕДРА\hspace{50pt} <<Теоретическая информатика и компьютерные технологии>>

\vspace*{-16pt}
\hspace{30pt}\rule{0.866\textwidth}{0.4pt}
  
\vspace{11em}

\begin{center}
\Large {\bf Летучка № 1} \\ 
\large {\bf по курсу <<Методы оптимизации>>} \\
\large <<Поиск минимума функции>> 
\end{center}\normalsize

\vspace{8em}


\begin{flushright}
  {Студент группы ИУ9-82Б Шемякин В.А. \hspace*{15pt}\\ 
  \vspace{2ex}
  Преподаватель Посевин Д. П.\hspace*{15pt}}
\end{flushright}

\bigskip

\vfill
 

\begin{center}
\textsl{Москва 2026}
\end{center}
\end{titlepage}
%--------------------------------------
%		КОНЕЦ ТИТУЛЬНОГО ЛИСТА
%--------------------------------------

\renewcommand{\ttdefault}{pcr}

\setlength{\tabcolsep}{3pt}
\newpage
\setcounter{page}{2}

\section{Задание}\label{Sect::task}
\begin{enumerate}
\item Изучить методы:
\begin{itemize}
\item метод Свенна локализации минимума унимодальной функции;
\item метод Фибоначчи поиска минимума унимодальной функции;
\item теоремы об условиях существования минимума функции.
\end{itemize}

\item Реализовать:
\begin{itemize}
\item Метод Свенна — локализация интервала, содержащего минимум.

\item Метод Фибоначчи — поиск минимума с заданной точностью. В консоль выводить и значение точки минимума, и заданную точность.

\item Численное доказательство унимодальности целевой функции вашего варианта на компакте, полученном методом Свенна. Первую производную находить только численно (аналитически нельзя). Целевая функция должна быть обёрнута в метод (функцию).

\item Численное доказательство (вторую производную аналитически находить нельзя), используя правило дождя, что найденная точка минимума действительно является минимумом, а не максимумом или точкой перегиба.
\end{itemize}
\end{enumerate}

\textbf{Вариант 29:} функция $f(x) = |x^3-3x|$.

\section{Результаты}\label{Sect::res}

Исходный код программы представлен в листинге 1.

\begin{breakablelisting}
function svenn(f, x0, h)
    x1 = x0
    x2 = x0 + h
    if f(x2) > f(x1)
        h = -h
        x2 = x0 + h
    end
    while f(x2) < f(x1)
        x1 = x2
        x2 = x1 + h
        h = 2 * h
    end
    a = min(x1, x2)
    b = max(x1, x2)
    return a, b
end

function fibonacci_search(f, a, b, eps)
    n = 1
    while fib(n + 2) < (b - a) / eps
        n += 1
    end
    k = 1
    x1 = a + (b - a) * fib(n) / fib(n + 2)
    x2 = a + (b - a) * fib(n + 1) / fib(n + 2)
    f1 = f(x1)
    f2 = f(x2)
    while k < n
        if f1 > f2
            a = x1
            x1 = x2
            f1 = f2
            x2 = a + (b - a) * fib(n - k + 1) / fib(n - k + 2)
            f2 = f(x2)
        else
            b = x2
            x2 = x1
            f2 = f1
            x1 = a + (b - a) * fib(n - k) / fib(n - k + 2)
            f1 = f(x1)
        end
        k += 1
    end
    candidates = [a, b, (a + b) / 2, x1, x2]
    x_min = candidates[1]
    for c in candidates
        if f(c) < f(x_min)
            x_min = c
        end
    end
    println("Точка минимума: ", x_min)
    println("Заданная точность: ", eps)
    return x_min, eps
end

function fib(n)
    if n <= 1
        return n
    end
    a, b = 0, 1
    for _ in 2:n
        a, b = b, a + b
    end
    return b
end

function num_derivative(f, x, delta=1e-7)
    return (f(x + delta) - f(x - delta)) / (2 * delta)
end

function num_second_derivative(f, x, delta=1e-6)
    return (f(x + delta) - 2 * f(x) + f(x - delta)) / (delta^2)
end

function prove_unimodality(f, a, b, n_points=100)
    step = (b - a) / (n_points - 1)
    sign_changes = 0
    prev_der = num_derivative(f, a)
    for i in 2:n_points
        x = a + (i - 1) * step
        der = num_derivative(f, x)
        if (prev_der < 0 && der >= 0) || (prev_der > 0 && der <= 0)
            sign_changes += 1
        end
        prev_der = der
    end
    unimodal = (sign_changes == 1)
    println("Численная проверка унимодальности на [", a, ", ", b, "]: смена знака первой производной ", sign_changes, " раз. Унимодальность: ", unimodal)
    return unimodal
end

function prove_minimum_rule(f, x_min, precision=1e-5)
    delta = max(1e-10, precision * 0.5)
    fpp = num_second_derivative(f, x_min, delta)
    f0 = f(x_min)
    rain_left = f(x_min - delta) >= f0 - 1e-12
    rain_right = f(x_min + delta) >= f0 - 1e-12
    is_min_rain = rain_left && rain_right
    is_min = is_min_rain
    println("Вторая производная в точке (численно): ", fpp)
    println("Правило дождя: f(x*-h) >= f(x*), f(x*+h) >= f(x*) при h=", delta, " : ", rain_left, ", ", rain_right)
    println("Точка является минимумом (не максимум и не перегиб): ", is_min)
    return is_min
end

f(x) = abs(x^3 - 3 * x)

a, b = svenn(f, 0.0, 0.1)
println("Метод Свенна: интервал [", a, ", ", b, "]")

x_min, precision = fibonacci_search(f, a, b, 1e-5)

prove_unimodality(f, a, b)

prove_minimum_rule(f, x_min, precision)

\end{breakablelisting}

\newpage

Результат представлен на рисунке~\ref{fig:img1}.

\begin{figure}[H]
  \centering
  \includegraphics[width=0.6\textwidth]{let1/screen1.png}
  \caption{Результат}
  \label{fig:img1}
\end{figure}.
\end{document}
